\documentclass[10pt]{article}

% ---------- Document Identity (update these only) ----------
\newcommand{\DocStage}{}
\newcommand{\DocTitle}{Your Road to Tier One \textbf{SOF}}
\newcommand{\ProgramName}{182-Week Civilian-to-Selection Readiness Pipeline}
\newcommand{\ProgramSubtitle}{}

% ---------- Page ----------
\usepackage[legalpaper,left=0.5in,right=0.5in,top=0.6in,bottom=0.45in]{geometry}

% ---------- Typography ----------
\usepackage[T1]{fontenc}
\usepackage[utf8]{inputenc}
\usepackage{libertinus}
\usepackage{microtype}
\usepackage{parskip}
\usepackage{ragged2e}
\usepackage{textcomp}
\AtBeginDocument{\color{black}}
\AtBeginDocument{\pagecolor{white}}
\AtBeginDocument{\justifying}

% ---------- Landscape pages ----------
\usepackage{pdflscape}

% ---------- Color ----------
\usepackage[table]{xcolor}
\definecolor{StageBlue}{HTML}{1F4E79}
\definecolor{StageBlueDark}{HTML}{163A5A}
\definecolor{LightGray}{HTML}{F2F4F7}
\definecolor{MidGray}{HTML}{D9DEE5}
\definecolor{RuleGray}{HTML}{C7CED8}
\definecolor{HeaderInk}{HTML}{000000}
\definecolor{Ink}{HTML}{000000}

% ---------- Utility ----------
\usepackage{enumitem}
\sloppy
\setlength{\emergencystretch}{2em}

% ---------- Boxes / UI ----------
\usepackage[most]{tcolorbox}

\tcbset{
  charterTitle/.style={
    enhanced,
    colback=StageBlue!4,
    colframe=StageBlue,
    boxrule=1.25pt,
    arc=3pt,
    left=16pt,right=16pt,top=14pt,bottom=14pt,
    borderline north={3.2pt}{0pt}{StageBlue},
    borderline west={4pt}{0pt}{StageBlue},
    borderline south={1.25pt}{0pt}{StageBlue},
    drop shadow={black!25!white},
  },
  charterRule/.style={
    enhanced,
    colback=LightGray,
    colframe=RuleGray,
    boxrule=0.85pt,
    arc=3pt,
    left=12pt,right=12pt,top=10pt,bottom=10pt,
    borderline west={3pt}{0pt}{StageBlue},
  },
  charterBadge/.style={
    enhanced,
    colback=StageBlue,
    coltext=white,
    colframe=StageBlueDark,
    boxrule=0.6pt,
    arc=2pt,
    left=6pt,right=6pt,top=3pt,bottom=3pt,
  }
}
\newtcbox{\Badge}{nobeforeafter,charterBadge}

% ---------- Lists ----------
\setlist[itemize]{leftmargin=1.5em, itemsep=0.25em, topsep=0.25em}
\setlist[enumerate]{leftmargin=1.8em, itemsep=0.25em, topsep=0.25em}

% ---------- TOC ----------
\usepackage{tocloft}
\setcounter{tocdepth}{3}
\setlength{\cftbeforesecskip}{3pt}
\setlength{\cftbeforesubsecskip}{2pt}
\setlength{\cftbeforesubsubsecskip}{0pt}
\renewcommand{\cftsecfont}{\bfseries}
\renewcommand{\cftsecpagefont}{\bfseries}
\renewcommand{\cftsubsecfont}{\normalfont}
\renewcommand{\cftsubsecpagefont}{\normalfont}
\renewcommand{\cftsubsubsecfont}{\normalfont}
\renewcommand{\cftsubsubsecpagefont}{\normalfont}
\setlength{\cftsecindent}{0pt}
\setlength{\cftsubsecindent}{1.25em}
\setlength{\cftsubsubsecindent}{2.8em}
\setlength{\cftsecnumwidth}{2.1em}
\setlength{\cftsubsecnumwidth}{2.8em}
\setlength{\cftsubsubsecnumwidth}{3.6em}

% Remove the built-in TOC heading so only the custom heading is shown
\makeatletter
\renewcommand{\tableofcontents}{\@starttoc{toc}}
\makeatother

% ---------- Header/Footer: page number only ----------
\usepackage{fancyhdr}
\pagestyle{fancy}
\fancyhf{}
\renewcommand{\headrulewidth}{0pt}
\renewcommand{\footrulewidth}{0pt}
\fancyfoot[C]{\thepage}

% ---------- Section formatting ----------
\usepackage{titlesec}

\titleformat{\section}
  {\Large\bfseries\color{Ink}}
  {\thesection}{0.6em}{}
  [\vspace{0.25em}\textcolor{StageBlue}{\rule{\linewidth}{0.8pt}}]

\titleformat{\subsection}
  {\large\bfseries\color{Ink}}
  {\thesubsection}{0.6em}{}

\titleformat{\subsubsection}
  {\normalsize\bfseries\color{Ink}}
  {\thesubsubsection}{0.6em}{}

\titlespacing*{\section}{0pt}{0.95em}{0.35em}
\titlespacing*{\subsection}{0pt}{0.65em}{0.25em}
\titlespacing*{\subsubsection}{0pt}{0.55em}{0.2em}

% ---------- Table engine ----------
\usepackage{tabularray}
\UseTblrLibrary{booktabs}

% ---------- tabularray: GLOBAL table treatment ----------
\SetTblrInner{
  cells = {fg=black, bg=white, valign=t},
}

% ---------- Header helpers ----------
\newcommand{\Hdr}[1]{\shortstack[c]{\strut #1\strut}}

% ---------- Lines ----------
\newcommand{\TitleLine}{\textcolor{StageBlue}{\rule{0.98\linewidth}{0.95pt}}}
\newcommand{\SubTitleLine}{\textcolor{RuleGray}{\rule{0.98\linewidth}{0.65pt}}}

% ---------- Continued on next page ----------
\newcommand{\ContinuedOnNextPageInline}{%
  \par
  \vfill
  \noindent\textcolor{RuleGray}{\rule{\linewidth}{1pt}}\vspace{-2pt}
  \noindent\hfill\textit{Continued on next page}\par\vspace{-10pt}
  \noindent\textcolor{RuleGray}{\rule{\linewidth}{1pt}}
  \clearpage
}

% ---------- PDF hyperlinks ----------
\usepackage[hidelinks]{hyperref}
\hypersetup{
  colorlinks=false,
  pdfborder={0 0 0},
  linktoc=all,
  pdftitle={\DocTitle},
  pdfauthor={\ProgramName}
}

% =========================================================
\begin{document}

\begin{tcolorbox}
\begin{center}
Stage 8 — Crosswalk Gap Check and Consistency Audit Packet\\[0.35em]
\textbf{SOF} Physical Development Transformation Program\\[0.25em]
182-Week Civilian-to-Selection Readiness Pipeline\\[0.65em]
\TitleLine
\end{center}
\end{tcolorbox}

\vspace{0.35em}

\section*{Stage 8 — Crosswalk Gap Check and Consistency Audit Packet}
\phantomsection
\addcontentsline{toc}{section}{Stage 8 — Crosswalk Gap Check and Consistency Audit Packet}

\subsection*{Introduction}
\phantomsection
\addcontentsline{toc}{subsection}{Introduction}

\subsubsection*{1) Purpose}

Stage 8 exists to perform system-level crosswalk QA across Stages 1–7 so the full package remains safe, auditable, and internally coherent. Stage 8 outputs are analysis, crosswalk mapping, and issue logging only. Stage 8 does not redesign the program and does not rewrite upstream artifacts.

Stage 8 is the mechanism that ensures gate decisions are defensible: the system must prove that the evidence required by governance can be produced using defined tests, inside defined windows, with defined resources available at time of need.

\subsubsection*{2) Scope}

\paragraph*{Inclusions}

\begin{itemize}
\item Stage 8 internal deliverable structure (listed for completeness; output below):
  \begin{itemize}
  \item 8.A Master Standards Index
  \item 8.B Standards to Tests to Phases Crosswalk
  \item 8.C Standards to Gear and Resources to Overviews Crosswalk
  \item 8.D Gap / Overload / Inconsistency Analysis Log
  \item 8.E Pre-Deploy and Implementation Checklists
  \end{itemize}
\item Issue logging using stable Stage 8 Issue IDs in the format \textbf{ISS-08}-\#\#\#.
\item Corrective principles limited to allowed moves only:
  \begin{itemize}
  \item adjust gate timing within Stage 3 protected windows,
  \item adjust tier timing within Stage 6,
  \item clarify mapping language in Stage 2, Stage 5, and Stage 7,
  \item do not rewrite protocols, thresholds, or phase purposes.
  \end{itemize}
\end{itemize}

\paragraph*{Exclusions}

\begin{itemize}
\item No new tests, thresholds, domains, phase purposes, or categories.
\item No rewriting of upstream artifacts.
\item No training programming content.
\end{itemize}

\subsubsection*{3) How to Use}

\paragraph*{Coach Lead}

\begin{itemize}
\item Run Stage 8 crosswalks before packaging/deployment and after any upstream patch that touches tests, gates, timing, or equipment readiness.
\item Treat Stage 8 outputs as a deployability decision layer:
  \begin{itemize}
  \item Must-fix issues (safety, inadmissible evidence, missing required resources) block deployment until resolved.
  \item Monitor issues (minor ambiguity, low-impact clarity gaps) may be carried forward with explicit documentation and constrained use posture.
  \end{itemize}
\item Use issue IDs to coordinate corrections and to prevent silent drift: any correction is traceable back to an \textbf{ISS-08}-\#\#\# entry.
\end{itemize}

\paragraph*{Stage 9 packaging}

\begin{itemize}
\item Stage 8 issues flow into Stage 9 QA summary and known issues reporting.
\item Stage 9 uses Stage 8 to declare coherence state: deployable vs draft-non-deployable, with explicit issue inventory.
\end{itemize}

\paragraph*{Execution}

\begin{itemize}
\item Stage 8 supports defensible gate decisions by verifying:
  \begin{itemize}
  \item evidence exists (\textbf{VALID} tests and admissible sessions),
  \item instruments exist (Stage 2 protocols and metric IDs),
  \item timing exists (Stage 3 protected windows),
  \item resources exist (Stage 6 Tier 1+ at time of need and Stage 7 operational manuals).
  \end{itemize}
\end{itemize}

\subsubsection*{4) Inputs and Assumptions}

\paragraph*{Inputs}

\begin{itemize}
\item Stage 1 governance and validity doctrine, including admissibility posture and gate outcomes.
\item Stage 2 authoritative measurement layer:
  \begin{itemize}
  \item domain taxonomy tags (\textbf{RUN}, \textbf{SWIM}, \textbf{CAL}, \textbf{RUCK}, \textbf{STR}, \textbf{BODYCOMP}, \textbf{DURABILITY}, \textbf{RECOVERY}),
  \item metric inventory with \textbf{MET-2B}-\#\#\#,
  \item protocol cards C\#,
  \item thresholds and logging conventions.
  \end{itemize}
\item Stage 3 calendars and deload/test windows and adjustment/reslot rules.
\item Stage 4 Phase 00 specification and exit standards.
\item Stage 5 macro design for Phases 01–13 and gate posture language.
\item Stage 6 tier timing matrices and phase-by-phase capability readiness.
\item Stage 7 category overviews 7.01–7.20.
\end{itemize}

\paragraph*{Assumption Ledger (only if needed)}

\begin{itemize}
\item Missing artifacts are recorded as gaps and logged as issues. No inference is permitted.
\end{itemize}

\subsubsection*{5) Interfaces}

\paragraph*{Upstream dependencies}

\begin{itemize}
\item Stage 8 audits Stages 1–7 for coherence, crosswalk completeness, timing alignment, and deployability readiness.
\end{itemize}

\paragraph*{Downstream consumers}

\begin{itemize}
\item Stage 9 QA summary, packaging doctrine, implementation guide, and deployment checklists.
\end{itemize}

\paragraph*{Crosswalk hooks}

\begin{itemize}
\item Stage 2 identifiers:
  \begin{itemize}
  \item Protocol cards C\#
  \item Metric IDs \textbf{MET-2B}-\#\#\#
  \end{itemize}
\item Domain tags:
  \begin{itemize}
  \item \textbf{RUN}, \textbf{SWIM}, \textbf{CAL}, \textbf{RUCK}, \textbf{STR}, \textbf{BODYCOMP}, \textbf{DURABILITY}, \textbf{RECOVERY}
  \end{itemize}
\item Phase identifiers and gate windows (Stage 3 protected windows).
\item Stage 7 category IDs 7.01–7.20 and Stage 6 tier timing (Tier 0–Tier 3, with Entry/Mid/End timing posture where used).
\end{itemize}

\subsubsection*{6) Gate and Acceptance Criteria for Stage 8}

Stage 8 is complete and deployable as an audit packet when:

\begin{itemize}
\item Crosswalk tables (8.A–8.C) are complete and every standard/metric is traceable to:
  \begin{itemize}
  \item at least one Stage 2 protocol card,
  \item at least one phase mapping that includes development and gate context,
  \item required resource categories at Tier 1+ available at time of need.
  \end{itemize}
\item Gap and conflict detection (8.D) is complete and all gaps/conflicts are logged using \textbf{ISS-08}-\#\#\# IDs with clear classification and disposition posture.
\item Corrective principles are provided at principle-level only and remain within allowed moves.
\item A must-fix vs monitor posture is produced so deployment decisions are explicit and auditable.
\end{itemize}

\subsubsection*{7) Common Failure Modes and Controls}

\begin{itemize}
\item Silent fixes instead of issue logging → prohibited; enforce \textbf{ISS-08}-\#\#\# logging for every detected mismatch and every correction.
\item Inventing tests or thresholds to patch gaps → prohibited; issues must be resolved by mapping, timing adjustments within allowed moves, or clarifying language.
\item Ignoring Stage 6 timing conflicts → prohibited; timing conflicts are logged and treated as deployability risks until resolved.
\item Missing join keys in crosswalk tables → enforce consistent use of C\#, \textbf{MET-2B}-\#\#\#, domain tags, phase IDs, and category IDs 7.01–7.20.
\item Understating safety conflicts → classify as Safety-type issues and treat as must-fix; safety overrides schedule and performance.
\item Confusing analysis with redesign → Stage 8 only flags and proposes corrective principles; does not rewrite upstream artifacts.
\end{itemize}

\end{document}